\documentclass[11pt,a4paper]{article}
\usepackage[margin=2.2cm]{geometry}
\usepackage{amsmath,amssymb,bm}
\usepackage{graphicx}
\usepackage{booktabs}
\usepackage{siunitx}
\usepackage[colorlinks=true,linkcolor=blue!50!black,citecolor=blue!50!black,urlcolor=blue!50!black]{hyperref}
\usepackage{xcolor}
\newcommand{\todo}[1]{\textcolor{red}{[TODO: #1]}}

\title{Composite A under finite-strain compression:\\
FFT-based FE simulation with true holes, material simplification,\\ and first steps towards third-medium contact}
\author{working notes (muFFTTO, branch \texttt{feat/fracture\_and\_contact})}
\date{29 October 2026}

\begin{document}
\maketitle

\begin{abstract}
We simulate the uniaxial compression of Composite~A of Li \& Rudykh~\cite{LiRudykh2019} -- a nearly incompressible
elastomer matrix with circular voids and stiff circular inclusions -- with the FFT-accelerated finite-element solver
muFFTTO (bilinear pixels with selective reduced integration, trust-region Newton--CG with the Green preconditioner).
With the voids modelled as true holes, the computed nominal stress--strain curve at 256 pixels per cell width matches
the experiment of~\cite{LiRudykh2019} up to $\varepsilon_X = 0.2$. We show that the material can be simplified --
inclusion/matrix stiffness contrast 10 instead of 1000 and matrix $\Lambda/\mu = 100$ instead of 1000 -- with a
change of the response of 1--3.5\,\% and a reduction of the solver cost by more than an order of magnitude. Filling
the voids with a third medium (for contact) raises the CG cost per Newton step by a factor 6--25; we trace this to the
mismatch between the bulk and the shear contrast of the compressible third medium and the nearly incompressible
matrix, and discuss remedies (augmented Lagrangian, mixed displacement--pressure formulation).
\end{abstract}

\tableofcontents

%=====================================================================================================================
\section{Problem}

\paragraph{Geometry.} The periodic primitive cell of Composite~A (\cite{LiRudykh2019}, Fig.~1a) has the size
$2a \times \sqrt{3}a$ with $a = 6.60$\,mm, i.e.\ $13.2 \times 11.4$\,mm. It contains three voids and one inclusion,
all of radius $r = 3$\,mm; the area fractions are 56.25\,\% voids, 18.75\,\% inclusion and 25\,\% matrix. The thinnest
matrix ligament is $0.6$\,mm. The 1\,$\times$\,1 primitive cell cannot show the buckling mode of Composite~A, which is
1\,$\times$\,2 (\cite{LiRudykh2019}, Sec.~4); we therefore compute a 1\,$\times$\,2 supercell
($13.2 \times 22.8$\,mm).

\paragraph{Loading.} Uniaxial compression in $X$: the macroscopic stretch $F_{XX} = 1 - \varepsilon$ is prescribed,
the lateral stretch $F_{YY}$ is an unknown of the energy minimization, which enforces $\langle P_{YY}\rangle = 0$
(lateral side stress free). The nominal stress is $-\langle P_{XX}\rangle$. The load is applied in increments
$\Delta\varepsilon = 0.005$ up to $\varepsilon = 0.2$; a random nodal imperfection of $10^{-3}$ pixel triggers the
buckling.

\paragraph{Material.} All phases are compressible neo-Hookean,
\begin{equation}
  W(\bm F) = \frac{\mu}{2}\,(I_1 - 2) - \mu \ln J + \frac{\lambda}{2}\,(\ln J)^2 ,
  \qquad I_1 = \operatorname{tr}(\bm F^{\mathsf T}\bm F),\; J = \det\bm F ,
  \label{eq:nh}
\end{equation}
in plane strain. Table~\ref{tab:material} lists the parameters of the paper and the simplified set used here
(Sec.~\ref{sec:simplification}). The voids are true holes (zero energy, stress and tangent).

\begin{table}[h]
\centering
\caption{Material parameters. $\nu = \lambda / (2(\lambda+\mu))$ is the plane-strain Poisson ratio.}
\label{tab:material}
\begin{tabular}{llrrrr}
\toprule
phase & & $\mu$ [MPa] & $\lambda$ [MPa] & $\lambda/\mu$ & $\nu$ \\
\midrule
matrix    & paper~\cite{LiRudykh2019} & 0.6 & 600 & 1000 & 0.4995 \\
          & simplified                & 0.6 & 60  & 100  & 0.495  \\
inclusion & paper (contrast 1000)     & 600 & 600\,000 & 1000 & 0.4995 \\
          & simplified (contrast 10)  & 6   & 600 & 100 & 0.495 \\
voids     & both                      & \multicolumn{4}{c}{true holes} \\
\bottomrule
\end{tabular}
\end{table}

%=====================================================================================================================
\section{Method}

\paragraph{Discretization.} The cell is discretized by a regular grid of bilinear (Q1) pixels with periodic
displacement fluctuation $\bm u$, $\bm F = \bar{\bm F} + \nabla\bm u$. Full Q1 integration locks for
$\Lambda/\mu = 1000$ (about 3$\times$ too stiff in our tests), so we use selective reduced integration (SRI): the
$\mu$-part of~\eqref{eq:nh} is integrated at the four Gauss points, the volumetric part $\frac{\lambda}{2}(\ln J_0)^2$
once per pixel with $J_0 = \det\bar{\bm F}_{\text{pixel}}$, the determinant of the pixel mean of $\bm F$. SRI gives the
same results as F-bar on the same grid and is more accurate per degree of freedom than Q2.

\paragraph{Nonlinear solver.} Equilibrium is the minimum of the total energy over $(\bm u, F_{YY})$. Each load
increment starts from the last converged state (no predictor) and is solved by trust-region Newton--CG
(NuMPI \texttt{tr\_newton\_bounded}, Steihaug CG). The Newton iteration stops when the infinity norm of the energy
gradient is below $10^{-4}\,\mu_m h\,\Delta\varepsilon$ ($h$: pixel size). A failed solve halves the load step.

\paragraph{Linear solver.} The Steihaug CG stops on the Hestenes--Stiefel lower estimate of the energy-norm error
with a fixed relative tolerance $10^{-6}$ and the adaptive delay of Meurant, Pape\v{z} \& Tich\'y with a minimum
delay of 10 (added to NuMPI for this study; it reproduces the stopping iteration of muFFTTO's CG exactly).
The preconditioner is the Green operator of the SRI Hessian at $\bm F = \bm I$ for the matrix material
(a translation-invariant operator inverted by FFT). All counts below are CG iterations with this setting.

\paragraph{Grids.} $N_X = 32, 64, 128, 256$ pixels per cell width ($32 \times 56$ up to $256 \times 444$ for the
supercell). At $N_X = 32$ the thinnest ligament spans only 1.4 pixels: the cell buckles at a stress 4$\times$ too low,
for any element. We use $N_X \ge 64$ for results; $N_X = 128, 256$ run on the GPU.

%=====================================================================================================================
\section{Material simplification}
\label{sec:simplification}

The paper's material combines two large ratios: the inclusion is 1000$\times$ stiffer than the matrix, and both are
nearly incompressible ($\Lambda/\mu = 1000$). Both slow down the CG. We tested whether they matter for the response.

\paragraph{Reasoning.}
\begin{itemize}
\item \emph{Inclusion contrast.} The inclusions are isolated (no continuous load path through the cell): the load
goes through the matrix ligaments around them. Once much stiffer than the matrix, an inclusion moves almost
rigidly; the effective stiffness saturates as the contrast grows, with a remainder of order $1/\text{contrast}$.
\item \emph{Matrix compressibility.} The cell deforms by bending and shearing of the thin ligaments, an isochoric
mode. $\lambda$ only has to be large enough to keep the volume nearly constant; it does not set the stiffness of the
structure, which is governed by $\mu$.
\item \emph{Inclusion bulk modulus.} We also set the inclusion bulk modulus equal to its shear modulus
($\lambda/\mu = 1/3$): no change of the response or the cost (not shown).
\end{itemize}

\paragraph{Results.} Table~\ref{tab:simplify128} and Fig.~\ref{fig:curves}a compare the simplified materials with
the paper's material at $N_X = 128$, where the paper's material is within about 15\,\% of~\cite{LiRudykh2019}.
Contrast 10 with $\Lambda/\mu = 100$ is 3.5\,\% softer before buckling, buckles at the same strain and is within 1\,\%
after buckling. $\Lambda/\mu = 10$ ($\nu = 0.45$) is 8--12\,\% too soft and changes the post-buckling shape (a stress
drop after buckling): the matrix is then noticeably compressible.

\begin{table}[h]
\centering
\caption{Nominal stress [MPa] at $N_X = 128$, and solver cost to $\varepsilon = 0.05$ (10 increments).
The paper-material run of 27 Oct used the earlier, looser inexact-Newton CG tolerance; its CG count is an
indication only.}
\label{tab:simplify128}
\begin{tabular}{lccc}
\toprule
$\varepsilon$ & paper material & contrast 10, $\Lambda/\mu = 100$ & contrast 10, $\Lambda/\mu = 10$ \\
\midrule
0.005 & 0.00172 & 0.00166 ($-3.5\,\%$) & 0.00151 ($-12\,\%$) \\
0.010 & 0.00341 & 0.00329 ($-3.5\,\%$) & 0.00300 ($-12\,\%$) \\
0.015 & 0.00464 & 0.00459 ($-1.1\,\%$) & 0.00445 ($-4\,\%$) \\
0.030 & 0.00489 & 0.00484 ($-1.0\,\%$) & 0.00447 ($-9\,\%$) \\
0.050 & 0.00524 & 0.00519 ($-1.0\,\%$) & 0.00480 ($-8\,\%$) \\
\midrule
CG to $\varepsilon = 0.05$ & $\approx 590\,000$ & 24\,665 & 8\,800 \\
GPU time & $\approx 7$\,h & 3.8\,min & 2.0\,min \\
\bottomrule
\end{tabular}
\end{table}

At $N_X = 32$ we also ran contrast $\{10, 100\}$ $\times$ $\Lambda/\mu$ $\{10, 100\}$ against the paper's material:
contrast 100 with $\Lambda/\mu = 100$ agrees with the paper's material to 2\,\% at 7$\times$ fewer CG iterations,
contrast 10 with $\Lambda/\mu = 100$ to 4\,\% at 32$\times$ fewer. (These $N_X = 32$ curves are themselves far from
the paper because of the under-resolved ligaments; they serve only the relative comparison.)

%=====================================================================================================================
\section{Results with true holes}

Figure~\ref{fig:curves}b shows the simplified material at $N_X = 256$ up to $\varepsilon = 0.2$
(40 increments, no failed Newton solve; 469 Newton steps, 211\,760 CG iterations, 1\,h\,52\,min on one GPU).

\begin{itemize}
\item The initial slope is 0.342\,MPa (paper simulation 0.35\,MPa).
\item Before buckling the curve lies within 3\,\% of the paper's material on the same grid
(Table~\ref{tab:256}).
\item After buckling it follows the \emph{experiment} of~\cite{LiRudykh2019} closely up to $\varepsilon = 0.2$
(0.0096 vs 0.0094\,MPa), below the paper's periodic simulation (0.0106\,MPa).
\item Buckling occurs late (overshoot to 0.0067\,MPa at $\varepsilon = 0.02$) because the imperfection is given in
pixels, i.e.\ it is 4$\times$ smaller in physical size at $N_X = 256$ than at 64. The paper-material run at 256 px
(other seed) buckled at 0.015. \todo{imperfection in physical units, rerun}
\item The deformation (Fig.~\ref{fig:snapshots}) shows the 1\,$\times$\,2 pattern: the voids collapse into curved
slits around the inclusions, alternating between neighbouring rows; compression concentrates in the ligaments. At
$\varepsilon \approx 0.2$ some slits are nearly closed: beyond this point true holes may overlap and contact is
required.
\end{itemize}

\begin{table}[h]
\centering
\caption{Nominal stress [MPa], $N_X = 256$.}
\label{tab:256}
\small
\begin{tabular}{lcccc}
\toprule
$\varepsilon$ & simplified, 256 px & paper material, 256 px & paper simulation & paper experiment \\
\midrule
0.005 & 0.00171 & 0.00177 & & \\
0.010 & 0.00339 & 0.00352 & & \\
0.015 & 0.00504 & 0.00513 & 0.0055 ($\varepsilon_{cr} = 0.016$) & \\
0.020 & 0.00665 (not yet buckled) & 0.00521 & 0.0054 & 0.0056 (peak at 0.018) \\
0.025 & 0.00525 & 0.00530 & & \\
0.050 & 0.00570 & & 0.0065 & 0.0051 \\
0.100 & 0.00673 & & 0.0077 & 0.0066 \\
0.150 & 0.00801 & & 0.0089 & 0.0079 \\
0.200 & 0.00960 & & 0.0106 & 0.0094 \\
\bottomrule
\end{tabular}
\end{table}

\begin{figure}[p]
\centering
\includegraphics[width=\textwidth]{figures/fig_curves.png}
\caption{Nominal stress--strain curves, compression in $X$, true holes. (a) Material simplification at
$N_X = 128$. (b) Simplified material (contrast 10, $\Lambda/\mu = 100$) at $N_X = 256$ against the paper's material
on the same grid (run to $\varepsilon = 0.025$) and against~\cite{LiRudykh2019}, Fig.~4.}
\label{fig:curves}
\end{figure}

\begin{figure}[p]
\centering
\includegraphics[width=\textwidth]{figures/fig_snapshots.png}
\caption{Deformed configuration, simplified material, $N_X = 256$. The computed 1\,$\times$\,2 cell (black box,
deformed by $\bar{\bm F}$) is drawn twice side by side. Matrix coloured by $P_{XX}$, inclusions grey, voids white.}
\label{fig:snapshots}
\end{figure}

\paragraph{Solver cost.} With true holes the Newton iteration takes 5--6 steps per increment before buckling,
28--43 in the buckling increment and 10--12 afterwards (Fig.~\ref{fig:cost}, left). Smaller load increments do not
reduce the per-increment count: the initial residual and the tolerance shrink together. The CG iterations per Newton
step grow from about 200 to 700 along the load path, almost identically at $N_X = 64$ and 256 (Fig.~\ref{fig:cost},
right): the Green preconditioner is grid-independent here.

%=====================================================================================================================
\section{Towards contact: the third medium}

Beyond $\varepsilon \approx 0.15$--$0.2$ the void faces touch. Third-medium contact fills the voids with a very soft
compressible neo-Hookean material ($\mu_{TM} = k\,\mu_m$, $\lambda_{TM}/\mu_{TM} = 1.5$, i.e.\ $\nu = 0.3$) and a
regularization $k_r(|\nabla\nabla\bm u|^2 - \tfrac12(\Delta\bm u)^2)$ inside the voids only ($k_r = 10^{-6} L^2 \mu_m$;
on Q1 pixels $\Delta\bm u$ vanishes, so only the Hessian term acts).

\paragraph{Effect on the response.} Before contact the medium only adds its own stiffness. With $k = 10^{-3}$ the
curve is 15--23\,\% above true holes after buckling and buckling is delayed: too stiff. With $k = 10^{-4}$ it is
within 3\,\% of true holes (64 px runs, to $\varepsilon = 0.03$ so far).

\paragraph{Effect on the cost.} The Newton iteration is unaffected (5--6 / 10--15 steps per increment), but the CG
iterations per Newton step rise from $\approx 450$ (true holes) to $\approx 2\,600$ ($k = 10^{-3}$) and
$\approx 4\,700$ ($k = 10^{-4}$), i.e.\ 6--25$\times$ (Fig.~\ref{fig:cost}).

\begin{figure}[h]
\centering
\includegraphics[width=\textwidth]{figures/contact_iterations.png}
\caption{Solver cost over the load path (Green preconditioner, simplified material): trust-region Newton steps per
increment, CG iterations per increment and CG iterations per Newton step, for true holes ($N_X = 64, 256$) and the
third medium $k = 10^{-3}, 10^{-4}$ ($N_X = 64$, still running).}
\label{fig:cost}
\end{figure}

\paragraph{Diagnosis: bulk/shear mismatch, not total contrast.} To isolate the linear solver we solved one Newton
system with the Hessian at $\bm F = \bm I$ (positive definite; at the compressed state the stiff inclusion makes the
Hessian indefinite by initial-stress softening) by plain PCG with the same stopping rule. Table~\ref{tab:linear}
collects the counts.

\begin{table}[h]
\centering
\caption{Linear test: CG iterations for $\bm H(\bm F = \bm I)\,\bm s = -\bm g$, Green preconditioner,
energy tolerance $10^{-6}$. ``$>$50k'': not converged within 50\,000 iterations.}
\label{tab:linear}
\begin{tabular}{llcc}
\toprule
material & third medium & $N_X = 32$ & $N_X = 64$ \\
\midrule
paper                         & true holes                                     & 1\,686 & 1\,323 \\
paper                         & $\mu_{TM} = 10^{-3}$, $\lambda_{TM} = 1.5\cdot10^{-3}$ & 28\,565 & 35\,322 \\
paper                         & $\mu_{TM} = 10^{-3}$, $\lambda_{TM} = 1$ ($\lambda/\mu$ as matrix) & 2\,299 & 2\,368 \\
paper                         & $\mu_{TM} = 1.5\cdot10^{-6}$, $\lambda_{TM} = 1.5\cdot10^{-3}$ & 43\,057 & \\
paper                         & $\mu_{TM} = 10^{-5}$, $\lambda_{TM} = 1.5\cdot10^{-5}$ & $>$50k & $>$50k \\
simplified (10, $\Lambda/\mu$ 100) & true holes                                & & 226 \\
simplified                    & $k = 10^{-3}$, $\nu = 0.3$                     & & 2\,661 \\
simplified                    & $k = 10^{-4}$, $\nu = 0.3$                     & & 6\,295 \\
simplified                    & $k = 10^{-4}$, Green reference $\lambda/\mu = 10$ & & 5\,756 \\
\bottomrule
\end{tabular}
\end{table}

With the same shear contrast ($\mu_{TM} = 10^{-3}$), raising only the bulk stiffness of the medium to the matrix ratio
cuts the count 12$\times$; lowering only its shear modulus makes it worse. The cost follows the \emph{worse} of the
two contrasts, measured against the reference medium in the corresponding mode:
\begin{equation}
  \frac{\mu_{TM}}{\mu_m} = k , \qquad
  \frac{\kappa_{TM}}{\kappa_m} = \frac{\lambda_{TM}+\mu_{TM}}{\lambda_m+\mu_m} \approx \frac{2.5\,k}{\Lambda/\mu} .
\end{equation}
For the nearly incompressible matrix the bulk contrast is $\Lambda/\mu$ times worse than the shear contrast. In the
simple contact benchmarks (patch, punch, Hertz) solids and medium had the same $\nu = 0.3$, so there was only one
contrast, which is why they were cheap. Moving the Green reference to an intermediate $\lambda/\mu$ does not help
(Table~\ref{tab:linear}, last row): an isotropic reference cannot match both phases. The inclusion's bulk modulus
plays no role.

\paragraph{Preconditioners tried.} Green-Jacobi fails inside the trust region at every load step size (its
$D^{-1/2}$ scaling spans $3\cdot10^{-2}$--$10^{6}$ and the trust region, measured in the preconditioner norm, keeps
the steps tiny), and is useless in the linear test ($>$50k). The dual preconditioner $G K(C_{\text{ref}}A^{-1}
C_{\text{ref}}) G$, for SRI exactly an SRI operator with the dual moduli $\mu_m^2/\mu$, $\lambda_m^2/\lambda$,
needs 3--4$\times$ fewer iterations than Green in the linear test (745 vs 2\,661 at $k = 10^{-3}$), but fails in the
trust region for the same reason as Green-Jacobi.

%=====================================================================================================================
\section{Literature and remedies}

We found no work combining third-medium contact with FFT/Green-type solvers. Third-medium contact is used with
finite elements and direct solvers~\cite{Bluhm2021,Wriggers2025}. For FFT methods, porous media~\cite{Schneider2020}
and pores combined with rigid inclusions (``infinitely double contrasted media'') are known to be hard, with a
convergence rate depending strongly on the reference medium. Near-incompressibility is addressed by
\begin{itemize}
\item the augmented Lagrangian of Michel, Moulinec \& Suquet~\cite{Michel2000} (fixed-point schemes);
\item the mixed deformation--pressure FFT-Galerkin formulation of Wang, Zhang \& Chen~\cite{Wang2022}: the pressure
is an unknown per voxel, the bulk modulus moves from the tangent into a diagonal block $-1/\kappa$, and the iteration
count stays flat (about 360) up to $\nu = 0.499$, where the standard formulation fails; hybrid incompressible
($\nu = 0.5$) / fully compressible ($\nu = 0$) foams are computed. The resulting system is a non-symmetric saddle
point problem, solved by CG without preconditioner and without theory; a preconditioner and inf--sup analysis are
left open.
\end{itemize}

\paragraph{Options for our problem.}
\begin{enumerate}
\item \emph{Augmented Lagrangian (next step).} Keep a pressure per pixel, but outside the CG: each inner problem is
an energy minimization with a moderate volumetric stiffness $r \ll \lambda_m$ (symmetric positive definite, our
trust-region Newton with Green unchanged, Green reference with $\lambda_{\text{ref}} = r$); an outer update of the
pressure restores the physical bulk modulus. Derivation in Appendix~\ref{app:al}.
\item \emph{Mixed displacement--pressure with a block preconditioner.} Keep the pressure in the linear system and use
MINRES with $\operatorname{diag}(G_{\mu}, (1/\mu + 1/\lambda)^{-1})$: Green with a shear-only reference for the
displacement block and a pointwise weight for the pressure block, robust in $\lambda/\mu$ in theory. Requires a
saddle-point nonlinear solver.
\item \emph{Material.} A less incompressible matrix ($\Lambda/\mu = 10$) reduces the mismatch 10$\times$ at an 8--12\,\%
error in the response.
\end{enumerate}

%=====================================================================================================================
\appendix
\section{Augmented Lagrangian for the volumetric term}
\label{app:al}

With $\theta = \ln J_0$ per pixel, the volumetric energy is the conjugate
$\frac{\lambda}{2}\theta^2 = \max_p \big(p\,\theta - \frac{p^2}{2\lambda}\big)$, $p = \lambda\theta$ the pressure.
The proximal-point (method of multipliers) iteration on $p$ with parameter $r$ gives, for a fixed $p_k$, the inner
energy
\begin{equation}
  \phi_k(\theta) = \max_p \Big( p\,\theta - \frac{p^2}{2\lambda} - \frac{(p - p_k)^2}{2r} \Big)
  = \frac{\kappa_r}{2}\Big(\theta + \frac{p_k}{r}\Big)^2 + \text{const}, \qquad
  \kappa_r = \frac{\lambda r}{\lambda + r},
\end{equation}
and the update $p_{k+1} = \phi_k'(\theta) = \kappa_r(\theta + p_k/r)$. Its fixed point is $p = \lambda\theta$, the
original problem. For $r \ll \lambda$ the inner volumetric stiffness is $\kappa_r \approx r$: each inner problem is a
well-conditioned energy minimization (the same SRI form with $\lambda \to \kappa_r$ and $\ln J_0 \to \ln J_0 + p_k/r$);
for a soft phase ($\lambda \ll r$, the third medium) $\kappa_r \approx \lambda$ and the phase is left essentially
unchanged. The outer iteration is the proximal point method on the concave dual and converges for any $r > 0$.

%=====================================================================================================================
\begin{thebibliography}{9}
\bibitem{LiRudykh2019} J.~Li, S.~Rudykh, Tunable microstructure transformations and auxetic behavior in 3D-printed
multiphase composites: the role of inclusion distribution, \emph{Composites Part B} 172 (2019) 352--362.
\bibitem{Wang2022} M.~Wang, K.~Zhang, C.~Chen, A mixed FFT-Galerkin approach for incompressible or slightly
compressible hyperelastic solids under finite deformation, \emph{Comput. Methods Appl. Mech. Engrg.} 396 (2022)
115092.
\bibitem{Michel2000} J.C.~Michel, H.~Moulinec, P.~Suquet, A computational method based on augmented Lagrangians and
fast Fourier transforms for composites with high contrast, \emph{CMES} 1 (2000) 79--88.
\bibitem{Schneider2020} M.~Schneider, Lippmann--Schwinger solvers for the computational homogenization of materials
with pores, \emph{Int. J. Numer. Methods Engrg.} 121 (2020) 5017--5041.
\bibitem{Bluhm2021} G.L.~Bluhm, O.~Sigmund, K.~Poulios, Internal contact modeling for finite strain topology
optimization, \emph{Comput. Mech.} 67 (2021) 1099--1114.
\bibitem{Wriggers2025} P.~Wriggers, J.~Kore\-lc, P.~Junker, A third medium approach for contact using first and
second order finite elements (2025).
\bibitem{Ladecky2025} M.~Ladeck\'y, I.~Pultarov\'a, F.~Bignonnet, I.~J\"odicke, J.~Zeman, L.~Pastewka,
Jacobi-accelerated FFT-based solver for smooth high-contrast data, arXiv:2508.02613 (2025).
\end{thebibliography}

\end{document}
